\documentclass[11pt]{article}
\usepackage{amsmath, amssymb, amsthm}
\usepackage{mathrsfs}
\usepackage{slashed}
\usepackage{hyperref}

\title{Operator Product Expansion:\\ A Revolution in Quantum Field Theory and the Renormalization Group}
\author{}
\date{}

\newtheorem{theorem}{Theorem}[section]
\newtheorem{definition}[theorem]{Definition}
\newtheorem{remark}[theorem]{Remark}

\begin{document}

\maketitle

\begin{abstract}
The Operator Product Expansion (OPE) provides a systematic framework for analyzing the short-distance behavior of products of local operators in quantum field theory. Introduced by Kenneth Wilson in the late 1960s, the OPE not only offered a rigorous foundation for the renormalization group but also transformed our conceptual understanding of quantum field theory itself. This self-contained introduction provides a detailed exposition of the OPE, including its definition, properties, and explicit construction in free field theory. It then elucidates the profound connection with the Wilsonian renormalization group and universality. Finally, it surveys major applications in conformal field theory, deep inelastic scattering, and effective field theory, demonstrating how the OPE reorganized the way we think about quantum fields.
\end{abstract}


\section{Introduction: The Challenge of Short Distances}
\label{sec:intro}

In quantum field theory (QFT), the fundamental dynamical variables are local operators $\mathcal{O}_i(x)$ constructed from the elementary fields and their derivatives. These operators represent physical observables: the electromagnetic current $J^\mu(x)$, the energy-momentum tensor $T^{\mu\nu}(x)$, or order parameters such as the local magnetization in a spin system. A central task in QFT is the computation of correlation functions of the form
\begin{equation}
\label{eq:correl}
\langle \mathcal{O}_1(x_1) \, \mathcal{O}_2(x_2) \cdots \mathcal{O}_n(x_n) \rangle .
\end{equation}
When two or more spacetime arguments approach each other, the correlation functions typically develop singularities. For instance, the two-point function of a free scalar field $\phi(x)$ of mass $m$ in $d$ Euclidean dimensions behaves as
\[
\langle \phi(x) \phi(0) \rangle = \frac{C_d}{|x|^{d-2}} + \text{regular as } x \to 0,
\]
with $C_d$ a known constant. In an interacting theory, the singularities are even more intricate, involving logarithms and operator mixing.

These short-distance singularities are not merely a mathematical nuisance; they encode essential ultraviolet (UV) information about the theory. In the pre-Wilsonian era, the dominant approach to handling them was renormalized perturbation theory: one introduced counterterms to subtract the infinities order by order in the coupling constant. While this prescription works remarkably well for quantum electrodynamics, it obscures the physical origin of scaling behavior and universality. Why do systems as disparate as a liquid-vapor critical point and a uniaxial ferromagnet share the same critical exponents? Why does the parton model describe deep inelastic scattering so successfully despite the confinement of quarks?

Kenneth Wilson's introduction of the operator product expansion (OPE) in 1969~\cite{Wilson:1969zs} provided a unified answer to these questions. The OPE asserts that the product of two local operators at nearby points can be expanded as a sum over a complete set of local operators with $c$-number coefficient functions that depend only on the separation. This deceptively simple statement has profound consequences. It elevates the short-distance singularity structure to an \emph{algebraic property} of the theory and provides the precise link between the UV and the infrared (IR) through the renormalization group.

In the following sections, we will develop the OPE in detail. We begin with its formal definition and elementary properties, then illustrate it in the solvable context of free field theory. Next, we explain Wilson's synthesis of the OPE and the renormalization group, showing how the expansion coefficients flow under scale transformations. Finally, we explore the wide-ranging impact of the OPE on conformal field theory, QCD phenomenology, and the modern philosophy of effective field theory.

\section{The Operator Product Expansion: Definition and Basic Properties}
\label{sec:definition}

\subsection{The Asymptotic Expansion}
Consider two local operators $\mathcal{O}_i(x)$ and $\mathcal{O}_j(y)$ in a quantum field theory. The OPE is the statement that, when inserted into any time-ordered (or radially-ordered, in Euclidean signature) correlation function, their product can be replaced by an asymptotic series as $x \to y$:
\begin{equation}
\boxed{
\mathcal{O}_i(x) \mathcal{O}_j(y) \underset{x \to y}{\sim} \sum_k C_{ij}^k(x-y) \, \mathcal{O}_k(y) .
}
\label{eq:OPEdef}
\end{equation}
Here:
\begin{itemize}
    \item The sum runs over a complete set of local operators $\{\mathcal{O}_k\}$ in the theory. By \emph{complete} we mean that any local operator can be expanded in this basis. The operators are usually chosen to be scaling operators with well-defined dimensions under scale transformations (in a scale-invariant theory) or eigenoperators of the renormalization group.
    \item The functions $C_{ij}^k(x-y)$ are ordinary (not operator-valued) functions of the separation vector. They are called \emph{Wilson coefficients}.
    \item The symbol $\sim$ indicates that this is an \emph{asymptotic expansion}. When the left-hand side is placed inside a correlation function $\langle \mathcal{O}_i(x) \mathcal{O}_j(y) \Phi(\infty) \rangle$, the right-hand side reproduces the exact correlation function up to terms that vanish faster than any power of $|x-y|$ (i.e., exponentially small or non-perturbative) as $x \to y$. In certain special theories like two-dimensional conformal field theories, the expansion actually converges at finite separation, but in general it is only asymptotic.
\end{itemize}

The OPE is an \emph{operator relation}. It does not depend on the specific correlation function in which it is inserted. This universality is its power: the same Wilson coefficients govern the short-distance behavior of the product in any Green's function, provided other operator insertions are kept far away compared to $|x-y|$.

\subsection{Scaling Dimensions and Power Counting}
If the theory possesses (perhaps approximate) scale invariance, the operators can be assigned scaling dimensions $\Delta_i$. Under a dilation $x \to \lambda x$, a local operator transforms as
\[
\mathcal{O}_i(\lambda x) = \lambda^{-\Delta_i} \mathcal{O}_i(x).
\]
In a free massless scalar theory, the field $\phi$ has dimension $\Delta_\phi = (d-2)/2$, while the composite operator $\phi^2$ has dimension $d-2$, etc. In an interacting theory, operators acquire \emph{anomalous dimensions} $\gamma_i$, so that their scaling dimension is $\Delta_i = \Delta_i^{(0)} + \gamma_i$, where $\Delta_i^{(0)}$ is the classical (engineering) dimension.

Dimensional analysis strongly constrains the form of the Wilson coefficients. Applying a dilation to both sides of Eq.~\eqref{eq:OPEdef} and requiring consistency, we find that $C_{ij}^k(x)$ must scale as
\begin{equation}
C_{ij}^k(\lambda x) = \lambda^{\Delta_k - \Delta_i - \Delta_j} \, C_{ij}^k(x).
\end{equation}
Therefore, the leading short-distance behavior is a power law:
\begin{equation}
C_{ij}^k(x) \sim |x|^{\Delta_k - \Delta_i - \Delta_j} \times (\text{possible logarithms}).
\label{eq:scaling}
\end{equation}
The exponent $\Delta_k - \Delta_i - \Delta_j$ is called the \emph{twist} in certain contexts. Operators with the smallest dimensions $\Delta_k$ give the most singular contributions as $x \to 0$. This provides a hierarchical organization: the OPE is dominated by the so-called \emph{relevant} and \emph{marginal} operators.

\subsection{Symmetry Constraints}
The OPE must respect all global symmetries of the theory. For example, if the theory has a conserved $U(1)$ current $J^\mu$, the product of two operators with definite charges will only contain operators whose charge is the sum of the individual charges. Likewise, Poincaré symmetry dictates that the Wilson coefficients can depend only on the Lorentz-invariant separation $|x-y|$ or, in the case of operators with spin, appropriate tensor structures constructed from $(x-y)^\mu$ and the metric. Conformal symmetry, when present, completely fixes the functional form of the Wilson coefficients up to a few constants (the OPE coefficients), as we shall see in Section~\ref{sec:CFT}.

\section{The OPE in Free Field Theory: A Solvable Example}
\label{sec:freefield}

To demystify the OPE, it is instructive to see how it arises in the simplest possible setting: the theory of a free massless scalar field in $d$ Euclidean dimensions. The action is
\[
S = \frac{1}{2} \int d^d x \, (\partial_\mu \phi)^2 .
\]
The fundamental two-point function is
\begin{equation}
\langle \phi(x) \phi(0) \rangle = \frac{\Gamma(d/2 - 1)}{4 \pi^{d/2}} \, \frac{1}{|x|^{d-2}} \equiv \frac{\mathcal{N}_d}{|x|^{d-2}} .
\label{eq:2pt}
\end{equation}
Consider the product of two fields $\phi(x)\phi(0)$. We can simply Taylor expand one of the fields around the origin:
\[
\phi(x) = \phi(0) + x^\mu \partial_\mu \phi(0) + \frac{1}{2} x^\mu x^\nu \partial_\mu \partial_\nu \phi(0) + \cdots .
\]
Multiplying by $\phi(0)$ gives
\[
\phi(x)\phi(0) = \phi^2(0) + x^\mu \phi(0)\partial_\mu \phi(0) + \frac{1}{2} x^\mu x^\nu \phi(0)\partial_\mu\partial_\nu \phi(0) + \cdots .
\]
This looks like an OPE, with the operators on the right being $\phi^2$, $\phi\partial_\mu\phi$, $\phi\partial_\mu\partial_\nu\phi$, etc. However, this naive expansion misses the singularities that arise when the product is inserted into a correlation function. Indeed, in free field theory, Wick's theorem tells us that
\[
\phi(x)\phi(0) = \langle \phi(x)\phi(0) \rangle \, \mathbf{1} + \, :\phi(x)\phi(0): \, ,
\]
where $\mathbf{1}$ is the identity operator and $:\phi(x)\phi(0):$ denotes normal ordering (subtraction of the self-contraction). The first term is the singular $c$-number part $\sim 1/|x|^{d-2}$. The second term is a well-defined composite operator. We can Taylor expand the normal-ordered product:
\[
:\phi(x)\phi(0): \, = \, :\phi^2(0): + x^\mu \, :\phi\partial_\mu\phi(0): + \frac{1}{2} x^\mu x^\nu \, :\phi\partial_\mu\partial_\nu\phi(0): + \cdots .
\]
Thus the exact OPE for the product $\phi(x)\phi(0)$ in free field theory is
\begin{equation}
\boxed{
\phi(x)\phi(0) = \frac{\mathcal{N}_d}{|x|^{d-2}} \, \mathbf{1} + \sum_{k} \frac{x^{\mu_1} \cdots x^{\mu_\ell}}{\ell!} \, :\phi \partial_{\mu_1}\cdots\partial_{\mu_\ell}\phi(0): .
}
\label{eq:freeOPE}
\end{equation}
The operators $:\phi \partial_{\mu_1}\cdots\partial_{\mu_\ell}\phi:$ have classical dimensions $\Delta_\ell = d - 2 + \ell$. According to the scaling rule~\eqref{eq:scaling}, the coefficient of such an operator should behave as $|x|^{\Delta_\ell - 2\Delta_\phi} = |x|^{\ell}$, precisely matching the explicit $x^{\mu_1}\cdots x^{\mu_\ell}$ factor.

This example illustrates several crucial points:
\begin{enumerate}
    \item The OPE contains a \emph{singular} $c$-number term multiplying the identity operator. This is the \emph{most singular} term because the identity has the smallest possible dimension $\Delta_{\mathbf{1}} = 0$.
    \item The sum over operators is an infinite series, but for small $|x|$, terms with larger $\ell$ (higher dimension) are suppressed by positive powers of $|x|$.
    \item The coefficients are completely determined by the kinematics (free field Wick contractions).
\end{enumerate}

\subsection*{A Composite Operator Example: $:\phi^2:(x) \, :\phi^2:(0)$}
To see a non-trivial OPE involving composite operators, consider the product $:\phi^2:(x) \, :\phi^2:(0)$. Using Wick's theorem, we contract fields from the first normal-ordered product with those from the second. There are two ways to make a single contraction (giving a term proportional to $:\phi(x)\phi(0):$) and one way to make two contractions (giving the square of the propagator times the identity). Expanding the resulting $:\phi(x)\phi(0):$ further yields:
\begin{equation}
\begin{aligned}
:\phi^2:(x) \, :\phi^2:(0) = & \; \frac{2\mathcal{N}_d^2}{|x|^{2d-4}} \, \mathbf{1} + \frac{4\mathcal{N}_d}{|x|^{d-2}} \, :\phi^2:(0) \\
& + \frac{4\mathcal{N}_d x^\mu}{|x|^{d-2}} \, :\phi\partial_\mu\phi:(0) + \cdots + \, :\phi^4:(0) + \cdots .
\end{aligned}
\end{equation}
Here the identity operator appears with a coefficient $\sim 1/|x|^{2d-4}$, which is more singular than in the $\phi\phi$ case. The operator $:\phi^2:$ itself appears with a coefficient $\sim 1/|x|^{d-2}$. This pattern of singularities is the free-field precursor to the mixing of operators under renormalization in interacting theories.

\section{Wilson's Renormalization Group and the OPE}
\label{sec:RG}

\subsection{The Wilsonian Effective Action}
Wilson's formulation of the renormalization group~\cite{Wilson:1971bg} is based on the idea of integrating out high-momentum degrees of freedom. Suppose we have a quantum field theory defined with an ultraviolet cutoff $\Lambda$ (e.g., a momentum cutoff or a lattice spacing). We wish to study the physics at a lower momentum scale $\mu \ll \Lambda$. To do so, we split the fields into ``fast'' modes with frequencies between $\mu$ and $\Lambda$, and ``slow'' modes with frequencies below $\mu$. Integrating out the fast modes yields an effective action for the slow modes, which can be expressed as a linear combination of local operators:
\begin{equation}
S_{\text{eff}}[\phi_{<}] = \sum_k g_k(\mu) \int d^d x \, \mathcal{O}_k(x) .
\label{eq:effaction}
\end{equation}
The coupling constants $g_k(\mu)$ are functions of the scale $\mu$. As we change $\mu$, the $g_k(\mu)$ flow according to the renormalization group equations. In particular, near a fixed point, the couplings scale as
\[
g_k(\mu) \sim \mu^{d - \Delta_k},
\]
where $\Delta_k$ is the scaling dimension of the operator $\mathcal{O}_k$. Operators with $\Delta_k < d$ are \emph{relevant} (their couplings grow in the infrared), those with $\Delta_k = d$ are \emph{marginal}, and those with $\Delta_k > d$ are \emph{irrelevant} (their couplings die out in the infrared).

\subsection{The OPE as the Bridge Between UV and IR}
How is the OPE connected to this picture? Consider a correlation function of two operators $\mathcal{O}_i(x)$ and $\mathcal{O}_j(0)$ with a set of other operators $\Phi$ located far away (at distances $\gg |x|$). When the separation $|x|$ is much smaller than the characteristic scale of the theory (e.g., the correlation length $\xi$), the physics at the scale $|x|$ is effectively decoupled from the IR physics. We can imagine performing the Wilsonian integration of modes with momenta between $1/|x|$ and $\Lambda$. The resulting effective theory at scale $1/|x|$ will contain an infinite number of local interactions. The effect of the short-distance product $\mathcal{O}_i(x)\mathcal{O}_j(0)$ on the far-away operators $\Phi$ is exactly captured by inserting the sum over local operators $\mathcal{O}_k(0)$ with coefficients given by the OPE coefficients $C_{ij}^k(x)$.

In equations, for any collection of far-away operators $\Phi$,
\begin{equation}
\langle \mathcal{O}_i(x) \mathcal{O}_j(0) \, \Phi \rangle = \sum_k C_{ij}^k(x) \, \langle \mathcal{O}_k(0) \, \Phi \rangle .
\label{eq:OPEcorrel}
\end{equation}
The right-hand side is precisely the evaluation of the correlation function in the effective theory at scale $\mu \sim 1/|x|$, where the couplings of the operators $\mathcal{O}_k$ have been renormalized to include the effect of the short-distance fluctuations. Therefore, the Wilson coefficients $C_{ij}^k(x)$ can be interpreted as the \emph{induced couplings} of the operators $\mathcal{O}_k$ when high-energy modes down to scale $1/|x|$ are integrated out.

This connection has profound consequences. The renormalization group equation for the Wilson coefficients follows from the independence of the physical correlation function on the arbitrary scale $\mu$ used to separate UV and IR. For a set of multiplicatively renormalizable operators, one finds
\begin{equation}
\left( \mu \frac{\partial}{\partial \mu} + \beta(g) \frac{\partial}{\partial g} + \gamma_i + \gamma_j - \gamma_k \right) C_{ij}^k(x;\mu,g) = 0,
\label{eq:RGE}
\end{equation}
where $\gamma_i$ are the anomalous dimensions of the operators. Solving this Callan--Symanzik-type equation resums the logarithms of $|x|$ that appear in perturbation theory, giving the correct scaling behavior $C_{ij}^k(x) \sim |x|^{\Delta_k - \Delta_i - \Delta_j}$.

\subsection{Universality Explained}
The OPE provides a natural explanation for \emph{universality} in critical phenomena. Near a second-order phase transition, the correlation length $\xi$ diverges, and the system becomes scale-invariant at distances much larger than the microscopic lattice spacing $a$ but much smaller than $\xi$. In this regime, the OPE is controlled by the fixed point of the renormalization group. The Wilson coefficients $C_{ij}^k(x)$ for the most relevant operators are \emph{universal}: they depend only on the fixed-point theory (the universality class) and not on the details of the microscopic Hamiltonian. Different physical systems that flow to the same fixed point will have identical leading OPE coefficients, and hence identical critical exponents and scaling functions. The subleading (irrelevant) operators give corrections to scaling that depend on the specific microscopic details, but these vanish as $\xi \to \infty$.

Thus, the OPE organizes the operators into a hierarchy of relevance, with the leading terms being universal and the subleading terms encoding non-universal corrections. This is why a liquid-gas critical point and a uniaxial ferromagnet can have the same critical exponents (they belong to the same Ising universality class), despite the completely different microscopic physics.

\section{Impact and Applications of the OPE}
\label{sec:applications}

The OPE, together with the Wilsonian RG, has fundamentally reshaped many areas of theoretical physics. We highlight three of the most significant arenas.

\subsection{Conformal Field Theory and the Bootstrap}
\label{sec:CFT}
In a conformal field theory (CFT), the OPE is not just an asymptotic short-distance expansion; it is an exact, convergent operator identity that holds at finite separation~\cite{Belavin:1984vu}. Conformal symmetry completely fixes the spacetime dependence of the Wilson coefficients. For scalar primary operators $\mathcal{O}_i$ of dimensions $\Delta_i$, the OPE takes the form
\begin{equation}
\mathcal{O}_i(x) \mathcal{O}_j(0) = \sum_k \frac{f_{ijk}}{|x|^{\Delta_i + \Delta_j - \Delta_k}} \left( \mathcal{O}_k(0) + \text{descendants} \right),
\end{equation}
where $f_{ijk}$ are the \emph{OPE coefficients} (constants) and the descendants (involving derivatives of $\mathcal{O}_k$) have completely determined coefficients. The entire dynamical content of a CFT is encoded in the spectrum of operator dimensions $\{\Delta_i\}$ and the OPE coefficients $\{f_{ijk}\}$.

This algebraic structure underlies the \emph{conformal bootstrap} program. By demanding that the OPE be associative---i.e., that the four-point function computed by applying the OPE in different channels (e.g., $12 \to 34$ versus $14 \to 23$) yields the same result---one obtains an infinite set of consistency conditions. These crossing symmetry equations severely constrain the allowed $\{\Delta_i, f_{ijk}\}$. In two dimensions, the infinite-dimensional Virasoro symmetry leads to exact solutions for many CFTs (the minimal models). In higher dimensions, numerical bootstrap techniques have produced world-record precision determinations of critical exponents in the 3D Ising model, all without relying on a Lagrangian description.

\subsection{Deep Inelastic Scattering and QCD}
\label{sec:QCD}
In Quantum Chromodynamics (QCD), the OPE is the rigorous field-theoretic foundation for the parton model and perturbative QCD predictions for hard processes. Consider deep inelastic scattering (DIS) of an electron off a proton, where a highly virtual photon with momentum $q$ ($Q^2 = -q^2 \gg \Lambda_{\text{QCD}}^2$) probes the proton structure. The hadronic tensor $W^{\mu\nu}$ is given by the Fourier transform of the proton matrix element of the product of two electromagnetic currents:
\[
W^{\mu\nu} \sim \int d^4 x \, e^{i q \cdot x} \langle P | J^\mu(x) J^\nu(0) | P \rangle .
\]
Because $Q^2$ is large, the dominant contribution comes from the light-cone region $x^2 \sim 1/Q^2 \to 0$. The OPE allows us to expand the product of currents in terms of local gauge-invariant operators:
\begin{equation}
J^\mu(x) J^\nu(0) \underset{x^2 \to 0}{\sim} \sum_{n} C_n^{\mu\nu}(x) \, \mathcal{O}_n(0),
\end{equation}
where the operators $\mathcal{O}_n$ are twist-two operators like $\bar{\psi} \gamma^{\{\mu_1} D^{\mu_2} \cdots D^{\mu_n\}} \psi$. The Wilson coefficients $C_n^{\mu\nu}(x)$ are singular as $x^2 \to 0$ and can be computed in perturbation theory thanks to asymptotic freedom. The proton matrix elements $\langle P | \mathcal{O}_n | P \rangle$ are non-perturbative quantities, but they are universal: they appear in many different processes and can be identified with the moments of parton distribution functions (PDFs). This \emph{factorization} of short-distance (Wilson coefficients) and long-distance (operator matrix elements) physics is the cornerstone of collider phenomenology.

The OPE also explains the existence of \emph{higher-twist} corrections (suppressed by powers of $1/Q^2$). These correspond to operators with larger twist $\tau = \Delta - s$ (dimension minus spin), whose contributions are subleading at high $Q^2$. The systematic inclusion of these corrections via the OPE is essential for precision QCD analyses.

\subsection{Effective Field Theory and Naturalness}
\label{sec:EFT}
The modern viewpoint of QFT is deeply rooted in the Wilsonian OPE. Any quantum field theory we write down should be regarded as an \emph{effective field theory} (EFT) valid up to some cutoff $\Lambda$. The Lagrangian contains all operators consistent with the symmetries, organized by their dimension:
\[
\mathcal{L}_{\text{eff}} = \mathcal{L}_{\text{renormalizable}} + \sum_{d > 4} \frac{c_i}{\Lambda^{d-4}} \mathcal{O}_i^{(d)} .
\]
The OPE provides the justification for this organization: operators of higher dimension are irrelevant in the IR, so their effects are suppressed by powers of $E/\Lambda$ at low energies $E \ll \Lambda$. This is why we can compute low-energy pion scattering using chiral perturbation theory without knowing the details of QCD at the GeV scale, and why we can analyze $B$-meson decays using the weak effective Hamiltonian without a complete theory of quantum gravity.

Furthermore, the OPE clarifies the concept of \emph{naturalness}. If a relevant operator (such as a scalar mass term) appears in the effective action, its coefficient is naturally of order $\Lambda^2$ (in units of the cutoff), unless protected by a symmetry. This is the essence of the hierarchy problem for the Higgs boson mass. The OPE tells us that the coefficient of the identity operator in the OPE of certain currents sets the scale for the mass term, and without a symmetry to forbid it, quantum corrections will drive it to the cutoff scale.

\section{Conclusion}
\label{sec:conclusion}

The Operator Product Expansion is far more than a technical tool for handling short-distance singularities. It encapsulates the deep physical principle that the microscopic details of a quantum field theory become encoded in a hierarchical expansion of local operators. By tying together the short-distance behavior of operators with the renormalization group flow of couplings, Wilson's OPE transformed renormalization from a set of ad hoc subtraction rules into a profound framework for understanding universality and scale dependence.

The OPE organizes our thinking about quantum field theory in an entirely new way. Instead of viewing a QFT as defined by a specific Lagrangian, we can regard it as defined by its operator content and the OPE coefficients among them. This algebraic perspective is at the heart of the conformal bootstrap, which has yielded unprecedented insights into strongly coupled systems. In particle physics, the OPE is the rigorous foundation for factorization theorems and precision QCD. In condensed matter, it explains why a myriad of microscopically distinct systems can exhibit identical critical behavior.

In summary, the Operator Product Expansion, in synergy with the renormalization group, has revolutionized the way we see quantum field theory. It shifted the focus from individual Feynman diagrams to the universal algebraic structure of local operators, providing a powerful language that unifies high-energy physics, critical phenomena, and even quantum gravity. Its legacy continues to shape the frontiers of theoretical physics.

\begin{thebibliography}{99}
\bibitem{Wilson:1969zs}
K.~G.~Wilson,
``Nonlagrangian models of current algebra,''
Phys.\ Rev.\ \textbf{179}, 1499-1512 (1969).

\bibitem{Wilson:1971bg}
K.~G.~Wilson,
``Renormalization group and critical phenomena. 1. Renormalization group and the Kadanoff scaling picture,''
Phys.\ Rev.\ B \textbf{4}, 3174-3183 (1971).

\bibitem{Belavin:1984vu}
A.~A.~Belavin, A.~M.~Polyakov and A.~B.~Zamolodchikov,
``Infinite Conformal Symmetry in Two-Dimensional Quantum Field Theory,''
Nucl.\ Phys.\ B \textbf{241}, 333-380 (1984).

\bibitem{Cardy:1996xt}
J.~L.~Cardy,
\textit{Scaling and Renormalization in Statistical Physics},
Cambridge University Press, 1996.

\bibitem{Collins:1984xc}
J.~C.~Collins,
\textit{Renormalization},
Cambridge University Press, 1984.

\bibitem{DiFrancesco:1997nk}
P.~Di Francesco, P.~Mathieu and D.~S\'en\'echal,
\textit{Conformal Field Theory},
Springer, 1997.
\end{thebibliography}

\end{document}